\section{Conclusion}
\label{sec:conclusion}

We introduced a geometric framework for understanding how optimal INR parameters move under signal group transformations, shifting the perspective from function-space equivariant constraints to parameter-space local tangent geometry.
Our key contributions are:

\begin{enumerate}
    \item \textbf{Low-dimensional orbital manifolds}: Signal group transformations induce low-dimensional manifolds in INR weight space, constrained by Lie algebra representation theory.
    \item \textbf{Perturbation stability bounds}: We quantified the symmetry error $\delta$ and demonstrated its empirical correlation with orbital regularity.
    \item \textbf{Tangent-aware fast adaptation}: A practical initialization strategy based on true orbital tangents achieves significant convergence speedup for conditional INR.
\end{enumerate}

The full arc of this work connects back to the puzzle that opened it.
COIN++ empirically achieved low-dimensional modulation representations without explaining why they suffice.
Our theory answers: because signal group transformations induce low-dimensional orbits in parameter space, and modulations are precisely coordinates along the orbit tangents.
This is not a coincidence of architecture or optimization---it is a geometric inevitability.

Looking forward, three directions hold particular promise.
First, extending the analysis to non-abelian Lie groups (e.g., $\SE(2)$ for combined rotation and translation) would reveal richer orbital topologies.
Second, geodesi-aware optimization that stays on the orbital manifold could replace gradient-based extrapolation entirely.
Third, applying this framework to dynamic scenes (NeRF with moving cameras) and video compression may unlock practical gains in representation efficiency.

\appendix

\section*{Appendix}
\label{sec:appendix}

\section*{Appendix A: Proofs}
\label{sec:appendix_theory}

\subsection*{Proof of Proposition 1 (Group Type Dictates Orbit Geometry)}

\begin{Proof}
Let $X \in \mathfrak{g}$ be the generator of a one-parameter subgroup $\exp(tX)$.
The tangent map $J_X$ satisfies the exponential map relation:
\begin{equation}
    \exp_*(tX) = e^{t J_X},
\end{equation}
where $\exp_*$ denotes the pushforward on the parameter manifold.

For $\mathfrak{g}_s \cong \mathbb{R}$ (scale group), $X$ has real eigenvalues $\{\lambda_i\}$.
Thus $e^{t\lambda_i}$ grows or decays monotonically in $t$, yielding open trajectories.

For $\mathfrak{so}(2) \cong \mathbb{R}$ (rotation group), $X$ has imaginary eigenvalues $\{\pm ik\}$.
Thus $e^{t(\pm ik)} = e^{\pm itk}$ oscillates periodically with period $2\pi/k$, yielding closed elliptical trajectories.
\qed
\end{Proof}

\subsection*{Proof of Perturbation Stability Lemma}

\begin{Proof}
Let $\tilde{J}_X = J_X + E$ where $\|E\| \le \delta \|J_X\|$.
By the stability theory of Kazhdan representations~\citep{kazhdan2012lovasz}, there exists an orthogonal $Q$ such that $Q^\top \tilde{J}_X Q$ is block-diagonal with error $O(\delta)$.
\qed
\end{Proof}

\section*{Appendix B: Implementation Details}
\label{sec:appendix_implementation}

\subsection{SIREN Architecture}

The SIREN network is defined as:
\begin{equation}
    f_\theta(\mathbf{x}) = W_3 \phi(W_2 \phi(W_1 \mathbf{x} + \mathbf{b}_1) + \mathbf{b}_2) + \mathbf{b}_3,
\end{equation}
where $\phi(x) = \sin(w_0 x)$ is the periodic activation with $w_0 = 30.0$.

\subsection{Jacobian-Vector Product Approximation}

The tangent vector is approximated via forward-mode automatic differentiation:
\begin{equation}
    \mathbf{v}_{\mathrm{tan}} = \frac{\partial \theta^*}{\partial s} \approx \frac{\partial \theta}{\partial s}\Big|_{\theta^*(s_1)},
\end{equation}
using PyTorch's \texttt{torch.autograd.functional.jvp}.

\subsection{MAML Configuration}

\begin{table}[h]
    \centering
    \caption{MAML hyperparameters used in experiments.}
    \label{tab:maml_config}
    \begin{tabular}{lc}
        \toprule
        Hyperparameter & Value \\
        \midrule
        Meta learning rate & $1e-3$ \\
        Inner learning rate & $5e-4$ \\
        Meta steps & 3000 \\
        Inner steps & 500 \\
        Fine-tune steps & 3000 \\
        \bottomrule
    \end{tabular}
\end{table}

\section*{Appendix C: Additional Experiments}
\label{sec:appendix_experiments}

\subsection{SIREN Width Ablation}

We vary the hidden dimension of SIREN and measure PC1 variance explained:

\begin{table}[h]
    \centering
    \caption{PC1 variance explained vs SIREN hidden dimension.}
    \label{tab:width_ablation}
    \begin{tabular}{cccc}
        \toprule
        Hidden Dim & Scale PC1 ($\%$) & Rotation PC1 ($\%$) & $\eps_{\mathrm{sym}}$ \\
        \midrule
        32 & \todo{result} & \todo{result} & \todo{result} \\
        64 & \todo{result} & \todo{result} & \todo{result} \\
        128 & \todo{result} & \todo{result} & \todo{result} \\
        \bottomrule
    \end{tabular}
\end{table}

\subsection{Seed Robustness}

We repeat experiments across 3 random seeds (42, 123, 456) and report standard deviations:

\begin{figure}[h]
    \centering
    \todo{placeholder for seed robustness plot}
    \caption{PC1 variance across random seeds. Error bars show $\pm 1$ standard deviation.}
    \label{fig:seed_robustness}
\end{figure}

